\section{A controlled core and a one-prime remainder}
\label{cg:section}

The central point of this section is a joint comparison. The core and
the remainder coordinates use one target hinge, even when many possible
remainder primes have large weights. The source restriction is that
each source uses at most one \emph{distinct} remainder prime; arbitrary
powers of that prime are permitted. No ordering or separation of prime
sizes is assumed.

All prime sets below are finite, all weights are nonnegative, and all
intervals have positive integer points. At real length $T>0$, a prefix
means $\{n\in\mathbb Z:1\le n\le T\}$ and an interval means the
integer points of $(x,x+T]$, $x\ge0$. For a prime set $F$, put
$F^\Pi=\prod_{q\in F}q$, with the empty product equal to one.

\subsection{Precise inherited analytic inputs}

We use the following conditional diffuse comparison from
\cite{Synthesis}. Its proof uses partial prime-power carriers and a
uniform interval version of the fundamental sieve lemma
\cite{Lagarias,GranvilleSound}. The statements below include the
quantifiers and constants needed here; the inherited proofs are included
in the companion archive.

\begin{proposition}[Conditional diffuse level comparison]
\label{cg:conditional-input}
There is an absolute constant $C>0$ for which the functions
\[
 e(u)=C\bigl((u/2)^{-u/4}+2^{1-u/2}\bigr)\quad(u\ge4),
 \qquad
 e_k(u)=3^ke(u/2)+3\cdot4^k2^{1-u/2}\quad(u\ge8)
\]
have the following property. Let $F$ contain at most $k$ primes, let
$L\cap F=\varnothing$, and put
\[
 V(n)=\sum_{p\in L}w_pv_p(n),\qquad
 0\le w_p\le\eta,\qquad
 \mu=\sum_{p\in L}\frac{w_p}{p-1}.
\]
For $\ell>\eta$, set
\[
 r=\lfloor\ell/\eta\rfloor\ge16,\qquad
 j=\lfloor r/2\rfloor,\qquad \kappa=r-j.
\]
If
\begin{equation}
 (1+e_k(\kappa))\mu+\eta2^{k+1-j}
 \le(\ell-\eta)(1-e_k(r)),
 \qquad e_k(r)<1,
 \label{cg:conditional-condition}
\end{equation}
then, for every $T>0$, every interval $J$ of real length $T$, and
every $s\ge0$,
\begin{equation}
 \#\{a\le T:(a,F^\Pi)=1, V(a)>s+\ell\}
 \le\#\{b\in J:(b,F^\Pi)=1, V(b)>s\}.
 \label{cg:conditional-level}
\end{equation}
The excluded primes may have arbitrary sizes, including sizes exceeding
$T$.

When $F=\varnothing$, the sharper unconditioned coefficient retained
in the same carrier proof is
\begin{equation}
 \#\{a\le T:V(a)>s+\ell\}
 \le
 \frac{(1+e(\kappa))\mu+\eta2^{1-j}}
      {(\ell-\eta)(1-e(r))}
 \#\{b\in J:V(b)>s\},
 \label{cg:unconditional-ratio}
\end{equation}
provided $e(r)<1$; here $r\ge8$ suffices.
\end{proposition}

For clarity, the numerator of the ratio in
\eqref{cg:unconditional-ratio}, divided by $\ell-\eta$, bounds
the source count per unit of a common carrier sum. The target count is
at least that carrier sum times $1-e(r)$. Thus the displayed quotient
is retained without a second loss in the comparison.

The following elementary alternative is also proved in
\cite{Synthesis}. Let $\mathbb E_0$ denote averaging over
$1,\ldots,M$. If the light support lies in $p\le M$, then
\begin{equation}
 \mu-3\eta\le\mathbb E_0V\le\mu,
 \qquad \mathbb E_0(V-\mu)^2\le8\eta\mu.
 \label{cg:moment-input}
\end{equation}
In particular, if $0<d\le1$, $\eta\le d^2/32$, and
$\mu\ge1+d$, then
\begin{equation}
 \sum_{a=1}^M\min(V(a),2)\ge M.
 \label{cg:high-mean-input}
\end{equation}
For any nonnegative raw load $X$, prime-power counts give
$\sum_I X\ge\sum_{a\le M}X(a)$. Consequently
\begin{equation}
 \sum_{a\le M}(X(a)-2)_+-\sum_{b\in I}(X(b)-1)_+
 \le M-\sum_{a\le M}\min(X(a),2).
 \label{cg:automatic}
\end{equation}
Thus \eqref{cg:high-mean-input} proves the full-prefix hinge inequality
at threshold two, and adding nonnegative coordinates preserves this
sufficient condition. Primes greater than $M$ must be removed before
computing the mean in this alternative. They change no source value;
restoring them only increases the target hinge.

\begin{lemma}[A threshold uniform in the number of exclusions]
\label{cg:uniform-eta}
For every $0<\varepsilon\le1$ there is $c_\varepsilon>0$ such that,
for every integer $k\ge0$, setting
\[
 \eta=\frac{c_\varepsilon}{k+1},\qquad \ell=1+\varepsilon
\]
makes \eqref{cg:conditional-level} valid whenever
$\mu\le1+\varepsilon/2$. We may also require
$c_\varepsilon\le\varepsilon^2/128$ for $0<\varepsilon\le1$.
\end{lemma}

\begin{proof}
For $u\ge16$,
\[
 e_k(u)\le3C3^k2^{-u/4}+6\cdot4^k2^{-u/2}.
\]
Indeed $(u/4)^{-u/8}\le2^{-u/4}$ for $u/4\ge4$.
Choose $R_\varepsilon\ge16$ so that
\[
 3C2^{-R_\varepsilon/4}+6\cdot2^{-R_\varepsilon/2}
 \le\varepsilon/16.
\]
Then $u\ge16k+R_\varepsilon$ implies
$e_k(u)\le\varepsilon/16$. Choose $c_\varepsilon$ with
\[
 c_\varepsilon\le\varepsilon^2/128,
 \qquad
 \frac{1+\varepsilon}{c_\varepsilon}
 \ge\max\{32,\,2R_\varepsilon+2\}.
\]
For $\eta=c_\varepsilon/(k+1)$, the integers in
Proposition~\ref{cg:conditional-input} satisfy
\[
 r\ge32k+2R_\varepsilon,\qquad
 \kappa\ge16k+R_\varepsilon,\qquad j\ge k+1.
\]
Hence both sieve errors are at most $\varepsilon/16$, and
$\eta2^{k+1-j}\le\eta$. At $\mu\le1+\varepsilon/2$, the
right-hand lower bound in \eqref{cg:conditional-condition} minus
its left-hand upper bound is at least
\[
 (1+\varepsilon-\eta)(1-\varepsilon/16)
 -(1+\varepsilon/16)(1+\varepsilon/2)-\eta
 \ge\frac{46}{128}\varepsilon-\frac{191}{2048}\varepsilon^2>0.
\]
Here $\eta\le\varepsilon/128$ was used. This proves the uniform
choice. Subsequent shrinking of $c_\varepsilon$ is harmless.
\end{proof}

\subsection{Gluing increments into one target hinge}

\begin{lemma}[One-prime source gluing]
\label{cg:gluing}
Let $R$ be a finite prime set, let $C(n)\ge0$ satisfy
$C(r^jt)=C(t)$ for every $r\in R$ and $j\ge0$, and let
$u_r:\mathbb Z_{\ge0}\to[0,1]$ be nondecreasing with $u_r(0)=0$.
Put $U(n)=\sum_{r\in R}u_r(v_r(n))$, $\ell>0$, and $c=1+\ell$.
Suppose $0<\rho\le1$ and $A_0,B_0\ge0$ have these properties:
\begin{align}
 \sum_{a\le M}(C(a)-c)_+
 &\le\rho\sum_{b\in I}(C(b)-1)_++A_0,
 \label{cg:base-assumption}\\
 \#\{a\le N:C(a)>s+\ell\}
 &\le\rho\#\{b\in J:C(b)>s\}+B_0
 \label{cg:level-assumption}
\end{align}
for all positive integer $N$, all $N$-term intervals $J$, and all
$0\le s\le1$. If $A\subseteq[2,M]$ is distinct and each source
has at most one distinct $R$-prime, then
\begin{equation}
 \sum_{a\in A}(C(a)+U(a)-c)_+
 \le\rho\sum_{b\in I}(C(b)+U(b)-1)_+
       +A_0+B_0|R|.
 \label{cg:gluing-conclusion}
\end{equation}
More precisely, the last error may be replaced by
$A_0+\sum_{r\in R}\min\{B_0,m_r\}$, where
$m_r=\#\{a\in A:r\mid a\}$.
\end{lemma}

\begin{proof}
For $h\in\{c,1\}$, define the nonnegative increment
\[
 D_{h,r}(n)=(C(n)+u_r(v_r(n))-h)_+-(C(n)-h)_+.
\]
Telescoping the profile and integrating the derivative of the hinge
gives the exact identity
\begin{equation}
 D_{h,r}(n)=\sum_{j\ge1}\mathbf1_{r^j\mid n}
  \int_{u_r(j-1)}^{u_r(j)}\mathbf1_{C(n)>h-v}\,dv.
 \label{cg:increment}
\end{equation}
For each $j$, divide multiples of $r^j$ by $r^j$. The source
quotients form a prefix of length $N=\lfloor M/r^j\rfloor$,
whereas the target quotients form an interval with at least $N$
consecutive positions. When $N>0$, take an $N$-term subinterval
and apply \eqref{cg:level-assumption} at $s=1-v\in[0,1]$.
The unused target positions contribute nonnegatively. When $N=0$,
the source count is zero. Invariance of $C$ under $r$-powers and
\eqref{cg:increment} now give
\begin{equation}
 \sum_{a\le M}D_{c,r}(a)
 \le\rho\sum_{b\in I}D_{1,r}(b)+B_0,
 \label{cg:per-prime}
\end{equation}
because $\sum_{j\ge1}(u_r(j)-u_r(j-1))\le1$.
Thus the error is paid once per prime, not once per prime power.

At each source, at most one remainder coordinate is nonzero, so
\begin{equation}
 (C(a)+U(a)-c)_+=(C(a)-c)_++\sum_{r\in R}D_{c,r}(a).
 \label{cg:source-identity}
\end{equation}
This identity is only needed on the source set. Every summand is
nonnegative, so distinctness permits us to enlarge source sums to the
prefix when applying \eqref{cg:base-assumption} and
\eqref{cg:per-prime}.

For $x,z_i\ge0$, one has
\begin{equation}
 (x-1)_++\sum_i\bigl((x+z_i-1)_+-(x-1)_+\bigr)
 \le\left(x+\sum_i z_i-1\right)_+.
 \label{cg:superadditive}
\end{equation}
If $x\ge1$ this is equality. Otherwise put $t=1-x>0$ and use
$\sum_i(z_i-t)_+\le(\sum_i z_i-t)_+$. Applying
\eqref{cg:superadditive} pointwise on the targets proves
\eqref{cg:gluing-conclusion}; the base target hinge has been used once.

For the refinement, the actual source increment
$S_r=\sum_{a\in A}D_{c,r}(a)$ is at most $m_r$, since each
increment is at most one and vanishes off the $r$-multiples.
It is also at most $\rho\sum_I D_{1,r}+B_0$ by
\eqref{cg:per-prime}. These two bounds imply
$S_r\le\rho\sum_I D_{1,r}+\min\{B_0,m_r\}$, proving the
sharper assertion.
\end{proof}

\begin{lemma}[Simultaneous restoration of raw tails]
\label{cg:tails}
Let $P_0$ be any finite set of coordinates capped at one, and write
\[
 Y(n)=V(n)+\sum_{p\in P_0}\min\{w_pv_p(n),1\},\qquad
 T(n)=\sum_{p\in P_0}(w_pv_p(n)-1)_+,
\]
where $V(n)\ge0$. If, for a distinct prefix subset $A$,
\[
 \sum_{a\in A}(Y(a)-c)_+\le\sum_{b\in I}(Y(b)-1)_++E_0,
\]
then the same estimate holds with $Y$ replaced by $Y+T$.
\end{lemma}

\begin{proof}
Pointwise, $(Y+T-c)_+\le T+(Y-c)_+$. Moreover,
\[
 (Y+T-1)_+=T+(Y-1)_+,
\]
because $T>0$ forces at least one capped coordinate to equal one,
and hence $Y\ge1$. Each summand of $T$ is a nondecreasing
prime-power profile vanishing at exponent zero. Expanding it in
nonnegative increments times $\mathbf1_{p^j\mid n}$ shows
\[
 \sum_{a\in A}T(a)\le\sum_{a\le M}T(a)\le\sum_{b\in I}T(b).
\]
Combining these inequalities proves the assertion. All tails are
restored in this single comparison.
\end{proof}

\subsection{Diffuse cores and one unrestricted core coordinate}

\begin{theorem}[A diffuse core with arbitrary one-prime remainders]
\label{thm:cg-diffuse}
For every $\varepsilon>0$ there is $\eta_\varepsilon>0$ with the
following property. Let $E,R$ be disjoint finite prime sets and put
\[
 X(n)=V(n)+W(n),\qquad
 V(n)=\sum_{p\in E}w_pv_p(n),\qquad
 W(n)=\sum_{r\in R}w_rv_r(n),
\]
where $0\le w_p\le\eta_\varepsilon$ for $p\in E$ and
$w_r\ge0$ is unrestricted for $r\in R$. If
$A\subseteq\{2,\ldots,M\}$ is distinct and each source is
divisible by at most one distinct $R$-prime, then, for every shift,
\begin{equation}
 \sum_{a\in A}(X(a)-2-\varepsilon)_+
 \le\sum_{b\in I}(X(b)-1)_+.
 \label{cg:diffuse-hinge}
\end{equation}
Arbitrary powers of that one remainder prime and arbitrary unweighted
cofactors are allowed. All thresholds are independent of the sizes and
cardinalities of the prime sets.
\end{theorem}

\begin{proof}
It suffices to treat $0<\varepsilon\le1$. Put
$\ell=1+\varepsilon$, $c=2+\varepsilon$, and
\[
 \rho_\varepsilon=\frac{1+3\varepsilon/4}{1+\varepsilon}<1.
\]
Choose $\eta=\eta_\varepsilon\le\varepsilon^2/128$ so small
that, with $r,j,\kappa$ as in
Proposition~\ref{cg:conditional-input},
\begin{equation}
 \frac{(1+e(\kappa))(1+\varepsilon/2)+\eta2^{1-j}}
      {(\ell-\eta)(1-e(r))}\le\rho_\varepsilon,
 \qquad e(r)<1,\quad r\ge8.
 \label{cg:strict-choice}
\end{equation}
Such a choice exists: the ratio tends to
$(1+\varepsilon/2)/(1+\varepsilon)$ as $\eta\downarrow0$.

Discard core primes greater than $M$, and compute
$\mu=\sum_{p\in E}w_p/(p-1)$ on the remaining support.
If $\mu\ge1+\varepsilon/2$, equations
\eqref{cg:high-mean-input} and \eqref{cg:automatic} give the
stronger full-prefix threshold-two inequality after adding $W$.
Restriction to $A$ proves the assertion.

Suppose $\mu\le1+\varepsilon/2$. Equations
\eqref{cg:unconditional-ratio} and \eqref{cg:strict-choice} give
\[
 \#\{a\le T:V(a)>s+\ell\}
 \le\rho_\varepsilon\#\{b\in J:V(b)>s\}
 \qquad(T>0, s\ge0).
\]
This comparison is uniform at every smaller length as well.
Integrating over $s\ge1$ gives its base hinge version at thresholds
$c$ and one. Apply Lemma~\ref{cg:gluing} with $C=V$,
$u_r(e)=\min\{w_re,1\}$, and $A_0=B_0=0$. This proves
\begin{equation}
 \sum_{a\in A}(V(a)+U(a)-c)_+
 \le\rho_\varepsilon\sum_{b\in I}(V(b)+U(b)-1)_+.
 \label{cg:strict-capped}
\end{equation}
Dropping the factor $\rho_\varepsilon\le1$ and applying
Lemma~\ref{cg:tails} to all $R$-coordinates restores their arbitrary
raw weights and powers. Finally restore the removed core primes on the
target side. This proves \eqref{cg:diffuse-hinge}.
\end{proof}

In the low-mean case, if every source has at most one remainder
\emph{occurrence} and $w_r\le1$, its remainder tail is zero.
Then \eqref{cg:strict-capped} yields the strict raw comparison with
factor $\rho_\varepsilon$ as well. Arbitrary remainder powers need
only the non-strict comparison in the theorem.

\begin{theorem}[One arbitrary core coordinate]
\label{thm:cg-one-heavy}
For every $\varepsilon>0$ there is $\eta_\varepsilon^*>0$ such
that Theorem~\ref{thm:cg-diffuse} remains valid when one prime
$q\in E$ has an arbitrary nonnegative raw weight and every other
core weight is at most $\eta_\varepsilon^*$. The number of
possible remainder primes and their weights are still unrestricted.
\end{theorem}

\begin{proof}
Again take $0<\varepsilon\le1$ and $\ell=1+\varepsilon$.
Write
\[
 V_0(n)=\sum_{p\in E\setminus\{q\}}w_pv_p(n),\qquad
 C(n)=V_0(n)+w_qv_q(n).
\]
Choose $\eta_\varepsilon^*$ no larger than the $k=1$ threshold
in Lemma~\ref{cg:uniform-eta}. Discard light primes above $M$
before computing their mean $\mu_0$.
The high-mean case $\mu_0\ge1+\varepsilon/2$ follows from
\eqref{cg:high-mean-input} and \eqref{cg:automatic} after adding
both the $q$-coordinate and the remainder load.

In the low-mean case, the conditional input gives
\begin{equation}
 \#\{a\le T:(a,q)=1, V_0(a)>s+\ell\}
 \le\#\{b\in J:(b,q)=1, V_0(b)>s\}
 \quad(s\ge0).
 \label{cg:one-conditional}
\end{equation}
We claim that the unrestricted core $C$ itself satisfies
\begin{equation}
 \#\{a\le T:C(a)>s+\ell\}
 \le\#\{b\in J:C(b)>s\}
 \quad(T>0, s\ge0).
 \label{cg:one-level}
\end{equation}
Decompose integers uniquely as $q^jt$ with $(t,q)=1$.
For every $j$ with $s-w_qj\ge0$, apply
\eqref{cg:one-conditional} at real length $T/q^j$, to the scaled
target interval and shifted level $s-w_qj$.
If $w_q=0$, these are all pieces. Otherwise the negative shifted
levels form the single suffix
\[
 j\ge j_0,\qquad j_0=\lfloor s/w_q\rfloor+1.
\]
On this suffix the target condition is automatic, since $V_0\ge0$.
The source count is at most $\lfloor T/q^{j_0}\rfloor$, whereas
the target count of the entire suffix is exactly the number of
$q^{j_0}$-multiples in $J$, which is at least that number.
All target exponents belong to this comparison, including $q^j>T$.
The boundary $w_qj=s$ belongs to the nonnegative-level case, not
to the suffix. Summing proves \eqref{cg:one-level} without error.

Integrating \eqref{cg:one-level} over $s\ge1$ gives the base
hinge comparison. The function $C$ is nonnegative and invariant
under every remainder-prime power. Lemma~\ref{cg:gluing}, now
with $\rho=1$ and zero errors, glues all capped remainder profiles.
Lemma~\ref{cg:tails} restores their tails. This proves the result.
The same suffix argument would permit any nondecreasing nonnegative
$q$-adic profile, although only the raw profile is needed here.
\end{proof}

\subsection{Finite variation of a capped exceptional upper set}

For several exceptional core coordinates, the negative-level pieces
need not form a single prime-power suffix. Their collective discrepancy
can nevertheless be bounded independently of the interval length.

\begin{lemma}[A finite divisor expansion]
\label{lem:cg-variation}
Let $F$ contain at most $k$ primes, each with $w_q\ge\eta>0$, and
put
\[
 Z(n)=\sum_{q\in F}\min\{w_qv_q(n),1\},\qquad
 f_s(n)=\mathbf1_{Z(n)>s}\quad(s\ge0).
\]
For every real $T>0$ and interval $J$ of length $T$,
\begin{equation}
 \left|\sum_{a\le T}f_s(a)-\sum_{b\in J}f_s(b)\right|
 \le4^k\binom{k+\lfloor s/\eta\rfloor}{k}.
 \label{cg:variation-bound}
\end{equation}
If $s\ge k$, or if $F$ is empty, the difference is zero.
\end{lemma}

\begin{proof}
The empty case is immediate. Write $F=\{q_1,\ldots,q_f\}$,
$f\le k$, and set $b_i=\lceil1/w_{q_i}\rceil$.
Let $u_i(j)=\min\{w_{q_i}j,1\}$ and
\[
 g_s(j_1,\ldots,j_f)=\mathbf1_{\sum_i u_i(j_i)>s}.
\]
For a nonnegative integer vector $\mathbf j$, define its mixed
divisor coefficient by
\begin{equation}
 a_{\mathbf j}(s)=
 \sum_{\substack{e_i\in\{0,1\}\\ e_i\le j_i\ (1\le i\le f)}}
 (-1)^{\sum_i e_i}g_s(\mathbf j-\mathbf e).
 \label{cg:finite-difference}
\end{equation}
Thus subtraction of one is allowed only in positive coordinates.
Telescoping separately in every coordinate gives the exact identity
\begin{equation}
 f_s(n)=\sum_{\mathbf j:\,\prod_iq_i^{j_i}\mid n}a_{\mathbf j}(s).
 \label{cg:divisor-expansion}
\end{equation}
If $j_i>b_i$ for any coordinate, the two adjacent layers of that
capped profile coincide, and its mixed difference is zero.
Consequently the expansion is finite. Also
$|a_{\mathbf j}(s)|\le2^f$.

For a nonzero vector $\mathbf j$, put $z_i=(j_i-1)_+$.
If $\sum_i u_i(z_i)>s$, all corners in
\eqref{cg:finite-difference} equal one, and their alternating sum
vanishes. The zero vector also has coefficient zero, since $s\ge0$
and $g_s(\mathbf0)=0$. A nonzero coefficient therefore requires
\[
 \sum_i u_i(z_i)\le s.
\]
On the relevant grid $j_i\le b_i$, one has $z_i<b_i$, so
$u_i(z_i)=w_{q_i}z_i\ge\eta z_i$. Hence every nonzero
coefficient satisfies
\begin{equation}
 \sum_i z_i\le s/\eta.
 \label{cg:corner-support}
\end{equation}
There are $\binom{f+\lfloor s/\eta\rfloor}{f}$ possible vectors
of nonnegative integers $z_i$ satisfying \eqref{cg:corner-support}.
Each zero coordinate $z_i=0$ permits $j_i=0$ or $1$; a positive
$z_i$ determines $j_i$ uniquely. Thus
\begin{equation}
 \sum_{\mathbf j}|a_{\mathbf j}(s)|
 \le4^f\binom{f+\lfloor s/\eta\rfloor}{f}
 \le4^k\binom{k+\lfloor s/\eta\rfloor}{k}.
 \label{cg:variation-norm}
\end{equation}

For each divisor $d$, its count in $J$ is either
$\lfloor T/d\rfloor$ or $\lfloor T/d\rfloor+1$.
Its prefix count is $\lfloor T/d\rfloor$. Each divisor-count
discrepancy thus has absolute value at most one. The triangle
inequality and \eqref{cg:variation-norm} prove
\eqref{cg:variation-bound} in both signs. This two-sided version
will be needed when the aggregate core tail is restored.
Finally, $Z(n)\le f\le k$, so $f_s$ vanishes for $s\ge k$.
\end{proof}

\subsection{Aggregating negative levels before estimating them}

\begin{lemma}[The collective exceptional-level defect]
\label{lem:cg-level-defect}
Let \(F,L\) be disjoint, let \(F\) contain at most \(k\) primes with
\(w_q\ge\eta>0\), and define
\[
 \begin{aligned}
 V(n)&=\sum_{p\in L}w_pv_p(n),&
 X_F(n)&=\sum_{q\in F}w_qv_q(n),\\
 Z(n)&=\min\{X_F(n),1\},& C(n)&=V(n)+Z(n).
 \end{aligned}
\]
Suppose the conditional comparison \eqref{cg:conditional-level}
holds for \(V,F,\ell\) at every nonnegative level and every real
length. Then, for all \(s\ge0\), all real \(T>0\), and intervals
\(J\) of length \(T\),
\begin{equation}
 \begin{split}
 &\#\{a\le T:C(a)>s+\ell\}-\#\{b\in J:C(b)>s\}\\
 &\qquad\le
 \sum_{a\le T}\mathbf1_{Z(a)>s}
       -\sum_{b\in J}\mathbf1_{Z(b)>s}.
 \end{split}
 \label{cg:aggregate-defect}
\end{equation}
In particular, the level defect for \(0\le s\le1\) is at most
\[
 B_{\eta,k}=4^k\binom{k+\lfloor1/\eta\rfloor}{k}
 \quad(k\ge1),
\]
and, with \(c=1+\ell\), the base hinge has no defect:
\begin{equation}
 \sum_{a\le M}(C(a)-c)_+-\sum_{b\in I}(C(b)-1)_+
 \le0.
 \label{cg:base-defect}
\end{equation}
The aggregate cap also satisfies
\begin{equation}
 \left|\sum_{a\le T}Z(a)-\sum_{b\in J}Z(b)\right|
 \le B_{\eta,k}.
 \label{cg:aggregate-cap-variation}
\end{equation}
For \(k=0\), \(B_{\eta,0}=0\).
\end{lemma}

\begin{proof}
Decompose each positive integer uniquely as \(n=ht\), where \(h\)
is its full \(F\)-smooth part and \((t,F^\Pi)=1\). Thus
\[
 C(ht)=Z(h)+V(t).
\]
For every \(h\) with \(Z(h)\le s\), apply the conditional input at
length \(T/h\), in the scaled target interval, at the nonnegative
level \(s-Z(h)\). The defect from this exact part is nonpositive.

For \(h\) with \(Z(h)>s\), the target condition holds automatically,
because \(V(t)\ge0\). The source contribution is at most the
unrestricted prefix count with this exact \(F\)-part. Sum all these
pieces before estimating their discrepancy. The sum is precisely
the difference of the two exceptional upper-set counts on the
right of \eqref{cg:aggregate-defect}. All target exact parts are
included, even those with \(h>T\). This proves the displayed
comparison; no error has been assigned separately to each smooth
part.

For \(0\le s<1\), the exceptional upper set is unchanged by
individual coordinate caps:
\[
 \mathbf1_{Z(n)>s}
 =\mathbf1_{X_F(n)>s}
 =\mathbf1_{\sum_{q\in F}\min\{w_qv_q(n),1\}>s}.
\]
Lemma~\ref{lem:cg-variation} therefore bounds its discrepancy in
both signs by \(B_{\eta,k}\). At \(s\ge1\), the upper set is
empty. This proves the level error on \(s\in[0,1]\).
Integrating \eqref{cg:aggregate-defect} over \(s\ge1\) gives
\eqref{cg:base-defect}, since its right side is now identically
zero. Finally,
\[
 Z(n)=\int_0^1\mathbf1_{Z(n)>s}\,ds.
\]
Integrating the two-sided discrepancy bound over this unit interval
proves \eqref{cg:aggregate-cap-variation}.
\end{proof}

\subsection{The uniform controlled-core theorem}

\begin{theorem}[Controlled substantial core coordinates]
\label{thm:cg-controlled}
For every \(\varepsilon>0\) there is \(c_\varepsilon>0\) such
that the following holds for every integer \(k\ge0\). Put
\[
 \eta=\frac{c_\varepsilon}{k+1}.
\]
Let \(E,R\) be disjoint finite prime sets. All but at most \(k\)
weights in \(E\) are at most \(\eta\); the remaining core weights
and all \(R\)-weights may be arbitrary nonnegative numbers. Define
\[
 X(n)=\sum_{p\in E\cup R}w_pv_p(n).
\]
If \(A\subseteq\{2,\ldots,M\}\) is distinct and each source is
divisible by at most one distinct \(R\)-prime, then
\begin{equation}
 \sum_{a\in A}(X(a)-2-\varepsilon)_+
       -\sum_{b\in I}(X(b)-1)_+
 \le B_{\varepsilon,k}(1+|R|)
 \label{cg:controlled-defect}
\end{equation}
for every positive \(M\)-term interval. Arbitrary powers of the
one remainder prime and arbitrary unweighted cofactors are allowed.
For \(k\ge1\), one may take
\begin{equation}
 B_{\varepsilon,k}
 =4^k\binom{k+\lfloor1/\eta\rfloor}{k}.
 \label{cg:error-constants}
\end{equation}
For \(k=0\), take \(B_{\varepsilon,0}=0\). The constant is
independent of \(M\), the prime values, the weights, and the interval
shift. One may restrict \(R\) to the primes occurring in the
sources. More precisely, the whole right side of
\eqref{cg:controlled-defect} may be replaced by
\begin{equation}
 B_{\varepsilon,k}+\sum_{r\in R}\min\{B_{\varepsilon,k},m_r\},
 \qquad m_r=\#\{a\in A:r\mid a\}.
 \label{cg:fibre-refinement}
\end{equation}
\end{theorem}

\begin{proof}
It suffices to treat \(0<\varepsilon\le1\). Choose
\(c_\varepsilon\) as in Lemma~\ref{cg:uniform-eta} and put
\[
 F=\{p\in E:w_p>\eta\},\qquad L=E\setminus F,\qquad
 V(n)=\sum_{p\in L}w_pv_p(n).
\]
Then \(|F|\le k\), and every \(F\)-weight is at least \(\eta\).
Discard light primes above \(M\) before computing
\(\mu_L=\sum_{p\in L}w_p/(p-1)\). If
\(\mu_L\ge1+\varepsilon/2\), the moment alternative
\eqref{cg:high-mean-input} applies because
\(\eta\le\varepsilon^2/128\). Adding all the remaining
nonnegative coordinates and using \eqref{cg:automatic} proves
\eqref{cg:controlled-defect} with zero error.

It remains to assume \(\mu_L\le1+\varepsilon/2\). The conditional
level comparison is now available uniformly at every real length,
with excluded set \(F\) and gap \(\ell=1+\varepsilon\).
Cap the whole substantial core at one, and cap each remainder
coordinate at one:
\[
 \begin{split}
 X_F(n)&=\sum_{q\in F}w_qv_q(n),\qquad
 Z(n)=\min\{X_F(n),1\},\qquad T_F(n)=X_F(n)-Z(n),\\
 U(n)&=\sum_{r\in R}\min\{w_rv_r(n),1\},\\
 T_R(n)&=\sum_{r\in R}(w_rv_r(n)-1)_+,\qquad
 C(n)=V(n)+Z(n).
 \end{split}
\]
Thus \(X=C+U+T_F+T_R\). Lemma~\ref{lem:cg-level-defect}
gives a level error \(B_{\varepsilon,k}\) on \(s\in[0,1]\)
and zero base hinge error. The function \(C\) is
invariant under every \(R\)-prime power. Lemma~\ref{cg:gluing},
with \(\rho=1\) and \(A_0=0\), therefore gives
\[
 \sum_{a\in A}(C(a)+U(a)-2-\varepsilon)_+
 \le\sum_{b\in I}(C(b)+U(b)-1)_+
       +B_{\varepsilon,k}|R|.
\]

The aggregate core tail is not a sum of separate prime-power
profiles, so it requires a separate justified estimate. Linear raw
prime-power counting and \eqref{cg:aggregate-cap-variation} give
\begin{equation}
 \begin{split}
 \sum_{a\le M}T_F(a)-\sum_{b\in I}T_F(b)
 &=\sum_{a\le M}X_F(a)-\sum_{b\in I}X_F(b)
       +\sum_{b\in I}Z(b)-\sum_{a\le M}Z(a)\\
 &\le B_{\varepsilon,k}.
 \end{split}
 \label{cg:aggregate-tail}
\end{equation}
Since \(T_F\ge0\), restriction to \(A\) preserves this upper
bound. The remainder tails retain exact domination,
\(\sum_A T_R\le\sum_I T_R\), as in Lemma~\ref{cg:tails}.
If \(T_F+T_R>0\), then \(C+U\ge1\): either \(Z=1\), or
one capped remainder coordinate equals one. Consequently
\[
 (X-1)_+=(T_F+T_R)+(C+U-1)_+,
 \qquad
 (X-2-\varepsilon)_+
 \le T_F+T_R+(C+U-2-\varepsilon)_+.
\]
Restoring these tails therefore costs at most one additional
\(B_{\varepsilon,k}\), not a separate cost for each exceptional
coordinate or smooth part. Restoring light primes above \(M\)
only increases the target load. This proves
\eqref{cg:controlled-defect}.

Primes of \(R\) absent from \(A\) may be removed before the proof
and restored only on the target side. The refined error
\eqref{cg:fibre-refinement} follows from the refined assertion of
Lemma~\ref{cg:gluing}; the aggregate tail restoration adds the
single \(B_{\varepsilon,k}\) in \eqref{cg:fibre-refinement}.
\end{proof}

The coefficient in the theorem is not asserted to be optimal.
In particular, Theorem~\ref{thm:cg-one-heavy} proves zero error
when \(k=1\) and the light threshold is chosen for that theorem.
The finite-variation estimate is designed to remain uniform for a
growing number of substantial core coordinates.

The error in \eqref{cg:aggregate-tail} cannot simply be replaced by
an assertion of exact prefix dominance for aggregate tails. Take
\(F=\{11,13,17\}\), all weights one, \(M=221\), and
\(I=\{2313,\ldots,2533\}\). The tail
\(T_F=(X_F-1)_+\) has prefix sum five: its positive contributions
occur at \(121,143,169,187,221\), each with value one. On \(I\),
the sum is four: \(2366\) and \(2420\) contribute one each, while
\(2431=11\cdot13\cdot17\) contributes two. This is a counterexample
to that intermediate dominance assertion, not to a hinge theorem
or to problem 708.

\begin{corollary}[Exponential growth of the error constant]
\label{cg:error-growth}
For \(k\ge1\), the constant above satisfies
\begin{equation}
 B_{\varepsilon,k}\le D_\varepsilon^k,\qquad
 D_\varepsilon=4e(1+2/c_\varepsilon).
 \label{cg:quantitative-errors}
\end{equation}
In the notation \(A_{\varepsilon,k}+B_{\varepsilon,k}|R|\),
one may take \(A_{\varepsilon,k}=B_{\varepsilon,k}\).
In particular, with no remainder primes, at most \(k\) substantial
weights cause a full-prefix hinge defect at most \(D_\varepsilon^k\),
uniformly in the ambient length.
\end{corollary}

\begin{proof}
For \(k\ge1\),
\[
 \lfloor1/\eta\rfloor\le2k/c_\varepsilon.
\]
Apply \(\binom Nk\le(eN/k)^k\) in
\eqref{cg:error-constants}. The no-remainder assertion follows
by taking the full distinct prefix as the source set in the theorem.
\end{proof}

\subsection{What the controlled estimate excludes}

Let \(n=|A|\to\infty\), and suppose the number of active remainder
primes satisfies \(|R|\le C_0n/\log n\), for a fixed constant
\(C_0\). Fix \(\varepsilon>0\) and \(0<\gamma<1\). If
\begin{equation}
 k\le\frac{(1-\gamma)\log\log n}{\log D_\varepsilon},
 \qquad
 \#\{p\in E:w_p>c_\varepsilon/(k+1)\}\le k,
 \label{cg:controlled-range}
\end{equation}
then
\[
 B_{\varepsilon,k}|R|
 =O\!\left(\frac{n}{(\log n)^\gamma}\right)=o(n),
 \qquad B_{\varepsilon,k}=O((\log n)^{1-\gamma})=o(n).
\]
Since \(x\le2+\varepsilon+(x-2-\varepsilon)_+\), the raw
dual objective of every such vector is at most
\begin{equation}
 \sum_{a\in A}X(a)-\sum_{b\in I}(X(b)-1)_+
 \le(2+\varepsilon)n+o(n).
 \label{cg:dual-exclusion}
\end{equation}
All error bounds here are independent of \(M\); no relation
between \(M\) and \(n\) is needed for this deduction. In the
minimal-obstruction class described later, the prime bound
\(p<cn\), with fixed \(c>2\), gives the required bound on
\(|R|\) by Chebyshev's estimate.

In particular, \eqref{cg:controlled-range} excludes a witness with
objective at least \((2+2\varepsilon)n\) for all sufficiently large
\(n\). It does not exclude an arbitrarily small excess over
\((2+\varepsilon)n\). This distinction is important when interpreting
an asymptotic estimate as a restriction on possible witnesses.
The theorem permits arbitrarily many substantial remainder weights;
the controlled count in \eqref{cg:controlled-range} is entirely
inside the designated core.

There is also an integer source-class consequence, proved with the
rounding step in Corollary~\ref{cor:tr-small-core}. If a designated
core has size \(o(\log\log M)\) and each source has at most one
distinct prime outside it, then
\(\tau(A,I)\le(2+o(1))|A|\), uniformly, without any separation of
prime sizes. In that application all core weights are unrestricted:
the number of substantial ones is bounded by the size of the core
for every dual vector.

That integer corollary records what follows from the controlled-weight
method. The separate construction in
Theorem~\ref{thm:sc-polylog} permits a designated core of size
\((\log M)^a\) for any fixed \(a<2\), and
Corollary~\ref{cor:sc-chebyshev-core} permits
\(c_0(\log M)^2/\log\log M\) core primes for an absolute \(c_0>0\).
Those constructions do not, by themselves,
increase the number of substantial coordinates handled uniformly
inside an unrestricted larger core in \eqref{cg:controlled-range}.

The unrestricted demand core can have size
\(O(\sqrt{n/\log n})\), much larger than the controlled range above.
The finite-variation errors therefore do not establish the
unrestricted asymptotic conjecture. A general proof would have to
bound their joint contribution more economically, or bypass these
upper-set discrepancies while retaining the single target budget.
